\section{A regularized matrix field}\label{sec:field}

We construct the field and primitive used in Section~\ref{sec:action}.
Throughout this section $K=\diag(k_1,\ldots,k_N)>0$ and $\delta>0$ are fixed.
The scalar case explains the choice of field. For $N=1$, write $K=(k)$.
If $-k\le b\le k$, putting $m=(k^2-b^2)^\sigma$ gives
\[
 m(k^2-b^2)=m^{1+1/\sigma}.
\]
Thus this field gives equality in the scalar version of the rank-one
comparison, and a primitive with derivative $-m$ has the required change
$\int_{-k}^k(k^2-b^2)^\sigma\,db$ from $k$ to $-k$.
We seek matrix versions of the comparison and endpoint identity,
with a regularization removed at the end.

The identity
\[
 k^2-2sb+s^2=(s-b)^2+k^2-b^2
\]
suggests integrating a scalar function of the quadratic on the left:
its affine dependence on $b$ will let us construct a trace primitive
even when the corresponding matrices $K$ and $B$ do not commute.
Set $\beta=\sigma-\tfrac12\in(-\tfrac12,\tfrac12)$ and define
$f:\R\to\R$ by
\begin{equation}\label{eq:f}
 f(0)=0,\qquad f'(y)=(\max\{y,0\}+\delta)^{\beta-1},
 \qquad
 c=\frac{\sigma}{\displaystyle\int_\R(1+s^2)^{\beta-1}\,ds}.
\end{equation}
The exponent and normalization are chosen so that, for $z\ge0$,
scaling the integration variable gives
\[
 c\int_\R f'(s^2+z)\,ds
 =\sigma(z+\delta)^{\beta-1/2}
 =\sigma(z+\delta)^{\sigma-1}.
\]
Integration in $z$ therefore produces the regularized power
$(z+\delta)^\sigma-\delta^\sigma$. On the scalar interval $[-k,k]$,
we will recover this expression with $z=k^2-b^2$.
Keeping $f'$ constant for negative arguments extends the construction
beyond that interval, as needed for the continuation argument.
The resulting $f$ is $C^1$, strictly increasing, and concave on $\R$.
The integral defining $c$ is finite because $\beta<1/2$; the definition
also includes $\beta=0$ without a limiting convention.

For $B\in\Sym_N$, put
\begin{equation}\label{eq:X-M}
 \begin{split}
 X_s(B)&=K^2-2sB+s^2I_N,\\
 M_R(B)&=c\int_{-R}^R\bigl(f(X_s(B))-f(s^2)I_N\bigr)\,ds,\\
 M(B)&=\lim_{R\to\infty}M_R(B).
 \end{split}
\end{equation}
Here and below scalar functions of symmetric or Hermitian matrices are
defined by the spectral theorem. The symmetric cutoff in
\eqref{eq:X-M} is part of the definition.

\subsection{Convergence and a primitive}

We first record an elementary estimate that keeps the regularity argument
in the Hilbert--Schmidt norm.

\begin{lemma}\label{lem:HS-calculus}
If a real function $h$ is $L$-Lipschitz on an interval containing the spectra
of two Hermitian matrices $A$ and $B$, then
$\|h(A)-h(B)\|_{\HS}\le L\|A-B\|_{\HS}$.
\end{lemma}
\begin{proof}
Choose orthonormal eigenbases $v_i$ and $w_j$ of $A$ and $B$, with respective
eigenvalues $\lambda_i$ and $\nu_j$. In these two bases the entries of
$h(A)-h(B)$ and $A-B$ are, respectively,
\[
 (h(\lambda_i)-h(\nu_j))\langle v_i,w_j\rangle,
 \qquad (\lambda_i-\nu_j)\langle v_i,w_j\rangle.
\]
Square, sum over $i,j$, and apply the scalar Lipschitz bound.
\end{proof}

\begin{lemma}\label{lem:field-regularity}
The limit in \eqref{eq:X-M} is locally uniform. The resulting map $M$ is
locally Lipschitz. Moreover,
\begin{equation}\label{eq:J}
 J(B)=-\int_0^1\tr\bigl(BM(tB)\bigr)\,dt
\end{equation}
is $C^1$ and satisfies \eqref{eq:framework-gradient}.
\end{lemma}
\begin{proof}
Fix a bounded set of $B$'s. For sufficiently large $|s|$, the spectrum of
$H_s(B)=K^2-2sB$ lies in $[-A|s|,A|s|]$, with $A$ independent of $B$.
On this interval set
\[
 h_s(t)=f(s^2+t)-f(s^2)-f'(s^2)t.
\]
All arguments of $f$ are positive. Since
$f''(y)=(\beta-1)(y+\delta)^{\beta-2}$ for $y>0$, the mean value theorem gives
\[
 \operatorname{Lip}(h_s)=O(|s|^{2\beta-3}),\qquad
 \sup_{|t|\le A|s|}|h_s(t)|=O(|s|^{2\beta-2}).
\]
It follows that
\begin{equation}\label{eq:tail}
 f(X_s(B))-f(s^2)I_N=-2s f'(s^2)B+R_s(B),
\end{equation}
where $R_s(B)=f'(s^2)K^2+h_s(H_s(B))$ has norm
$O(|s|^{2\beta-2})$. By Lemma~\ref{lem:HS-calculus}, its Lipschitz constant
as a function of $B$ has the same bound: the factor $2|s|$ from $H_s$ multiplies
$O(|s|^{2\beta-3})$. The first term on the right of \eqref{eq:tail} is odd
in $s$ and cancels at every symmetric cutoff; the remainder is integrable.
On bounded $s$-intervals, $f$ is Lipschitz and the same lemma applies directly.
This proves both assertions about $M$.

For completeness we justify the primitive without differentiating $M$.
Let $F$ be any $C^2$ primitive of $f$. The trace differentiation formula is
\begin{equation}\label{eq:trace-derivative}
 d\tr F(A)[H]=\tr(f(A)H).
\end{equation}
For polynomials it follows by cyclicity of trace. Approximation of $F$ and
$F'$ on a compact spectral interval by a polynomial and its derivative
proves the general formula. For $s\ne0$, the map
$B\mapsto f(X_s(B))-f(s^2)I_N$ is therefore a gradient: a potential is
\[
 -\frac{1}{2s}\tr F(X_s(B))-f(s^2)\tr B.
\]
At $s=0$ the field is constant and is again a gradient. Consequently each
$M_R$ has zero integral around every piecewise $C^1$ closed curve, by
integration first along the curve and then in $s$. Its negative has the
primitive
$J_R(B)=-\int_0^1\tr(BM_R(tB))\,dt$.
Local uniform convergence gives $J_R\to J$ and $M_R\to M$ on compact sets.
Passing to the limit in the line-segment identity for $J_R$ yields
\[
 J(B+H)-J(B)=-\int_0^1\tr(M(B+tH)H)\,dt.
\]
Continuity of $M$ proves that $J$ is $C^1$ with the stated derivative.
\end{proof}

\subsection{Scalar comparison and boundary values}

Define
\begin{equation}\label{eq:mu}
 \mu(z)=c\int_\R\bigl(f(s^2+z)-f(s^2)\bigr)\,ds,
 \qquad z\in\R.
\end{equation}
This integral is absolutely convergent: on bounded $z$-sets its integrand
has tail $O(|s|^{2\beta-2})$.

\begin{lemma}\label{lem:mu}
The function $\mu$ is locally Lipschitz and strictly increasing, with
$\mu(0)=0$. For $z\ge0$,
\begin{equation}\label{eq:mu-positive}
 \mu(z)=(z+\delta)^\sigma-\delta^\sigma,
 \qquad \mu'(z)=\sigma(z+\delta)^{\sigma-1}.
\end{equation}
For every fixed $b,z\in\R$,
\begin{equation}\label{eq:shift}
 c\lim_{R\to\infty}\int_{-R}^R
 \bigl(f((s-b)^2+z)-f(s^2)\bigr)\,ds=\mu(z).
\end{equation}
\end{lemma}
\begin{proof}
The tail estimate above, also for differences in $z$, proves local
Lipschitz continuity. Strict increase follows from strict increase of $f$
at every $s$. Differentiation under the integral for $z\ge0$ gives
\[
 \mu'(z)=c\int_\R(s^2+z+\delta)^{\beta-1}\,ds
 =\sigma(z+\delta)^{\beta-1/2}
 =\sigma(z+\delta)^{\sigma-1}.
\]
Integration from zero proves \eqref{eq:mu-positive}.

For the shift, write $g(s)=f(s^2+z)$. It is even, and
$g'(s)=2s f'(s^2+z)=O(|s|^{2\beta-1})\to0$ as $|s|\to\infty$.
Suppose first that $b\ge0$. The change in the integral of $g$ under the
shift is
\[
 \int_R^{R+b}g(t)\,dt-\int_{R-b}^{R}g(t)\,dt.
\]
Its absolute value is at most
$b^2\sup_{|t-R|\le b}|g'(t)|$, which tends to zero. Evenness handles
$b<0$ as well. Subtracting $f(s^2)$ now proves \eqref{eq:shift}.
\end{proof}

\begin{lemma}[Directional comparison]\label{lem:jensen}
For every $B\in\Sym_N$ and every real unit vector $e$,
\begin{equation}\label{eq:jensen}
 e^{\transpose}M(B)e\le
 \mu\bigl(e^{\transpose}K^2e-(e^{\transpose}Be)^2\bigr).
\end{equation}
If a standard basis vector $e_i$ is an eigenvector of $B$ with eigenvalue
$b_i$, then
\begin{equation}\label{eq:common-eigenvector}
 M(B)e_i=\mu(k_i^2-b_i^2)e_i.
\end{equation}
For the primitive in \eqref{eq:J},
\begin{equation}\label{eq:boundary-J}
 J(-K)-J(K)=\sum_{i=1}^N\int_{-k_i}^{k_i}\mu(k_i^2-s^2)\,ds.
\end{equation}
\end{lemma}
\begin{proof}
Scalar concavity and the spectral decomposition of $X_s(B)$ give
\[
 e^{\transpose}f(X_s(B))e\le f(e^{\transpose}X_s(B)e).
\]
Set $b=e^{\transpose}Be$ and $z=e^{\transpose}K^2e-b^2$. The scalar argument
on the right is $(s-b)^2+z$. Integrate at a symmetric cutoff and use
Lemma~\ref{lem:mu} to obtain \eqref{eq:jensen}. This uses scalar concavity
only; no operator concavity on the whole real line is asserted.

For \eqref{eq:common-eigenvector}, $e_i$ is also an eigenvector of $X_s(B)$,
with eigenvalue $(s-b_i)^2+k_i^2-b_i^2$. Functional calculus and
\eqref{eq:shift} give the identity. Finally, restrict $dJ=-\tr(M\,dB)$ to
diagonal matrices and integrate along the diagonal segment from $K$ to
$-K$. Each coordinate contributes
\[
 -\int_{k_i}^{-k_i}\mu(k_i^2-s^2)\,ds
 =\int_{-k_i}^{k_i}\mu(k_i^2-s^2)\,ds,
\]
as claimed.
\end{proof}

In particular $M(K)=M(-K)=0$. We shall also use the following symmetry.
If $S$ is a diagonal matrix with entries $\pm1$, then $SKS=K$ and
\begin{equation}\label{eq:sign-equivariance}
 M(SBS)=SM(B)S.
\end{equation}
This follows already at each cutoff. The matrix $K$ stays fixed throughout;
we do not require equivariance under conjugations that move $K$.
